\section{Smaller local graphs for the direct checker}
\label{sec:direct-local-wave11}

The latest pinned repository supplies a complete numeric checker for the
original sealed-ZFC calculus, independent of K215 and Jones finalization.
Its $83{,}281$ raw and $76{,}313$ optimized primitive-recursive program nodes
are not polynomial variable counts.  This section improves concrete local
polynomial graphs used by that route, without changing its numeric codes.
Every visible variable and witness below ranges over $\mathbb N$, including
zero; every subtraction is ordinary integer-polynomial subtraction.

\begin{center}
\begin{tabular}{lrrr}
\hline
Local graph & Old witnesses/ops & New witnesses/ops & Degree\\
\hline
Max-square pair & $3/46$ & $2/39$ & $4$\\
Pair, strict squared rows & $3/46$ & $2/42$ & $4$\\
Either unpair projection & $4/46$ & $3/39$ & $4$\\
Predecessor & $0/16$ & $0/10$ & $4$\\
Tree left/right, list tail & $5/62$ & $4/49$ & $4$\\
If-zero conditional & $2/32$ & $1/16$ & $4$\\
Zero test & $1/20$ & $1/14$ & $4$\\
Actual list head & $7/76$ & $5/59$ & $4$\\
Pair, selector eliminated & $3/46$ (quartic) & $1/53$ & $6$\\
\hline
\end{tabular}
\end{center}
Operations count binary additions, subtractions and multiplications in the
literal expression tree, including the final addition of zero, as in the
baseline.  Repeated subtrees are counted repeatedly.  The artifacts also
provide separate shared-DAG counts and full expanded integer coefficients.
No local saving is advertised as a completed smaller consistency polynomial.

\subsection{Remove the maximum-component witness}
For natural $z,a,b,t,s$, put
\[
 B=t(t-1),\qquad L=a-b-s+t(2s+1),\qquad
 E=4(z-a-b)-(a+b+s)^2+t,
\]
\[
 \operatorname{Pair}(z,a,b;t,s)=B+L^2+E^2.
\]
Then
\[
 \exists t,s\in\mathbb N:\operatorname{Pair}=0
 \quad\Longleftrightarrow\quad
 z=\begin{cases}b^2+a,&a<b,\\a^2+a+b,&b\le a.\end{cases}
\]
Indeed $t(t-1)\ge0$ for every integer $t$, so a zero forces $t=0$ or
$t=1$ and $L=E=0$.  At $t=0$, $a=b+s$ and $a+b+s=2a$; hence
$z=a^2+a+b$.  At $t=1$, $b=a+s+1$ and $a+b+s=2b-1$; hence
$z=b^2+a$.  Conversely these two branches choose respectively
$(t,s)=(0,a-b)$ and $(1,b-a-1)$.  This covers ties and zero.
The literal quartic has two witnesses, $39$ tree operations and $52$
monomials.  Replacing $B$ by $B^2$ gives a strict sum of three squared
quadratic rows, with $42$ operations and the same zero set.  This latter
version fits the existing squared-row interface without a nonnegative-term
extension.  Hiding the missing component yields either unpair projection
with three witnesses.

Solving the last row for the selector also gives the one-witness sextic
\[
 T=(a+b+s)^2-4(z-a-b),\qquad
 T(T-1)+\bigl(a-b-s+T(2s+1)\bigr)^2=0.
\]
Here $T$ is an integer polynomial, not a new coordinate.  Its nonnegative
first summand forces $T\in\{0,1\}$; the preceding branch proof then applies.
This has $53$ tree operations and $86$ monomials, but raises degree to six.

\subsection{Share the predecessor and remove branch selectors}
The exact predecessor graph factors as
\[
 [n(p+1-n)]^2+[p(p+1-n)]^2=(n-p-1)^2(n^2+p^2).
\]
Its zero set is precisely $p=\max(n-1,0)$: either $n=p+1$ or $n=p=0$.
Adding this product to $\operatorname{Pair}(p,a,b;t,s)$ gives either tree
projection with four hidden coordinates and $49$ operations.  At $n=0$
both components are zero, so the right projection is also the actual guarded
list tail.

A separate selector is unnecessary for a conditional.  The quartic
\[
 [n(c-b)]^2+[(c-a)(n-s-1)]^2=0
\]
with one natural witness $s$ is equivalent to $c=a$ if $n=0$, and $c=b$
otherwise.  At zero, $n-s-1=-s-1\ne0$; at positive $n$, the first row
forces $c=b$, and completeness chooses $s=n-1$.  Taking $a=1,b=0$ gives
the $14$-operation zero test.  If an exact predecessor $p$ is already
available, set $s=p$ and allocate no extra slack.

For the actual found-head convention, with visible $n,c$ and witnesses
$p,a,b,t,s$, use
\[
 \operatorname{Head}=\operatorname{Pair}(p,a,b;t,s)
 +(n-p-1)^2(n^2+p^2+c^2)+[n(c-a-1)]^2.
\]
At $n=0$, its middle term forces $p=c=0$ and the pairing term forces
$a=b=0$.  At $n>0$, it forces $p=n-1$, $c=a+1$, and the correct unpair
components.  Conversely these choices make every term zero.  Thus nil is
zero and a found zero-valued entry is one.  The ordinary quartic has five
witnesses, $59$ operations and $72$ monomials.

\subsection{Certificate status}
The package supplies literal trees and full coefficient lists for eleven
local polynomials.  Exact symbolic identities, counts and $1{,}845{,}448$
finite controls pass in ordinary and optimized Python modes; checks remain
active under \texttt{-O}.  Universal correctness follows from the arguments
above, not from the finite controls.  These new proofs have not been checked
by Lean, and the whole numeric checker, its bounded loops and its internal
nonstandard-code interpretation still require integration.  Natural witness
domains are essential; unrestricted integer witnesses need their own
naturalness encoding.

\section{Index-free inert-prime closure with a reused product scale}
\label{sec:index-free-crt-wave11}

Fix $m\ge1$, natural $n$, external natural parameters $x$, and fixed
nonnegative-coefficient polynomials $e(i,w;x),f(i,w;x)$.  Set $P=e-f$ and
choose a fixed $d\ge0$ bounding its degree in $i$; for $P=0$, $d=0$ is
allowed.  For a natural witness bound $U$, define
\[
\begin{split}
 H&=e(n,U,\ldots,U;x)+f(n,U,\ldots,U;x)+n+U+1,\\
 T&=H!,\qquad D=4T,\qquad S=D^n n!,\qquad M=DS+1,\\
 B&=\bigl(S\tbinom{S+n}{n}\bigr)\bmod M.
\end{split}
\]
Writing $P(i,w;x)=\sum_{k=0}^dP_k(w;x)i^k$, set
\[
 R_D(A;x)=\sum_{k=0}^dP_k(A;x)D^{d-k}(1-D)^k.
\]
This is an ordinary polynomial.  Its inverse-looking shorthand
$D^dP(D^{-1}-1,A;x)$ introduces no actual inverse or variable exponent.
For a fixed nested binary norm $N$, obtained by repeatedly replacing a pair
of subexpressions by the sum of their squares, the exact certificate is
\[
 \exists U,A_1,\ldots,A_m\in\mathbb N:\quad
 B\mid N\left(\tbinom{A_1}{U+1},\ldots,
                 \tbinom{A_m}{U+1},R_D(A;x)\right).
\]
It is equivalent to $\forall i<n\ \exists w\in\mathbb N^m:P(i,w;x)=0$.
Binomial is total, with $\binom Ak=0$ for $A<k$.  Neither factorials nor
binomials nor $D^n$ are being called ordinary polynomial notation.

\paragraph{Why the product needs only one variable-power graph.}
For $m_i=D(i+1)-1$,
\[
 S\binom{S+n}{n}=\prod_{j=1}^n(DS+Dj)
 \equiv\prod_{j=1}^n(Dj-1)\pmod{DS+1}.
\]
For $n>0$ the positive product is smaller than $D^n n!=S<M$; at $n=0$
it is one, still below $M$.  Hence $B=\prod_{i<n}m_i$.  The former
independent bound $(Dn)^n+1$ is unnecessary: reuse $S$ itself.  This removes
one complete variable-power graph, not merely its displayed abbreviation.

\paragraph{Why the encoded index is unnecessary.}
The integer polynomial $R_D-D^dP(i,A;x)$ is divisible by $D(i+1)-1$,
by the difference-of-powers identity.  Thus modulo $m_i$, the cleared
residual is exactly $D^dP(i,A;x)$.  The modulus already encodes the index
as $D^{-1}-1$, so no CRT coordinate $c$ or index norm leaf is needed.

Each $m_i$ is $3$ modulo $4$, all its prime factors exceed $H$, and the
$m_i$ are pairwise coprime: a common divisor divides $j-i$, hence $H!$,
but is coprime to $D$.  Choose a prime $p_i\equiv3\pmod4$ dividing $m_i$.
Anisotropy of a binary sum of squares over $\mathbb F_{p_i}$ makes every
norm leaf zero.  The binomial leaves recover residues $w_j\in[0,U]$;
$D$ is invertible, so the last leaf gives $P(i,w;x)=0\pmod{p_i}$.
Its original absolute value is smaller than $H<p_i$, hence it is zero
as an integer.  The bound is applied after cancelling $D^d$, not to the
large cleared residual.  Conversely CRT representatives of actual witnesses
make every leaf divisible by every $m_i$, proving completeness.  The empty
bound $n=0$ is handled by $B=1$.

\paragraph{An exposed, conditional witness ledger.}
Suppose chosen factorial, power and total-binomial graph polynomials use
respectively $\nu_F,\nu_P,\nu_B$ internal witnesses beyond visible arguments
and output.  Expose graph outputs
\[
 T=H!,\quad F=n!,\quad E=D^n,\quad C=\binom{EF+n}{n},\quad
 L_j=\binom{A_j}{U+1}.
\]
Use ordinary expressions $S=EF$, $M=DS+1$, and natural $B,K,r,V$ satisfying
\[
 SC=B+MK,\qquad M=B+r+1,\qquad N(L_1,\ldots,L_m,R_D)=BV.
\]
There are exactly $2m+9$ exposed coordinates.  The complete conditional
count is
\[
 2m+9+2\nu_F+\nu_P+(m+1)\nu_B,
\]
versus $2m+11+2\nu_F+2\nu_P+(m+1)\nu_B$ for the preceding interface
under the same convention.  The saving is exactly $\nu_P+2$, but no
numerical values for these uncompiled graph costs have been asserted.
This is a proved compiler simplification, not yet a smaller closed ordinary
consistency polynomial.  Further quartic flattening and internal formalization
have separate costs and proof obligations.
